\documentclass[11pt]{article}

% ---- Packages (all standard in a full TeX Live / MiKTeX install) ----
\usepackage[utf8]{inputenc}
\usepackage[T1]{fontenc}
\usepackage{lmodern}
\usepackage[margin=1in]{geometry}
\usepackage{amsmath,amssymb,amsthm}
\usepackage{booktabs}
\usepackage{array}
\usepackage{multirow}
\usepackage{microtype}
\usepackage{xcolor}
\usepackage{listings}
\usepackage{siunitx}
\usepackage[hidelinks]{hyperref}
\usepackage[capitalize,noabbrev]{cleveref}

\sisetup{group-separator={,},group-minimum-digits=4,detect-weight=true}

% ---- Macros ----
\newcommand{\bt}[1]{\texttt{#1}}
\newcommand{\Z}{\mathbb{Z}}
\newcommand{\Fabb}{\mathcal{F}_{\mathsf{ABB}}}
\newcommand{\ReLU}{\ensuremath{\mathrm{ReLU}}}
\newcommand{\Baru}{\textsc{Baru}}
% A finding worth pulling out of the prose.
\newcommand{\finding}[1]{\medskip\noindent\fbox{\parbox{0.97\linewidth}{\textbf{Finding.} #1}}\medskip}

\lstset{basicstyle=\ttfamily\small,columns=fullflexible,keepspaces=true,
        breaklines=true,frame=none,xleftmargin=1em}

\title{\bfseries Measuring Privacy-Preserving Computation in One Place:\\
MPC, FHE, ZKP and Their Hybrids}

\author{(author names elided for review)}
\date{}

\begin{document}
\maketitle

\begin{abstract}
Privacy-preserving computation is not one field but four---secure multi-party
computation (MPC), fully homomorphic encryption (FHE), garbled circuits, and
zero-knowledge proofs (ZKP)---each with its own literature, benchmarks and
accounting conventions. A practitioner choosing between them, or a reviewer
assessing a claim that one beats another, must assemble numbers from different
papers, measured on different hardware, over different networks, counting
different things. Cross-paradigm comparison is therefore largely \emph{citation
assembly} rather than measurement.

We take a different route. We survey the paradigms through a single artifact,
\Baru{}, in which one source program runs on all of them---a plain interpreter,
three-party replicated-secret-sharing MPC, garbled circuits, BFV/BGV/TFHE
homomorphic encryption, a Binius64 SNARK, and \emph{combinations} in which
values cross between representations mid-computation without being exposed.
Every backend is driven through the same transport, charged by the same byte and
round accounting, and checked bit-exactly against the same plaintext reference.
A network cost model is validated against a real socket to within \SI{13.6}{\percent}
worst case, so a claim about a wide-area deployment is computed from measured
counts rather than asserted.

Several widely repeated intuitions survive re-measurement; several do not. Among
the results that only a single-harness study can produce: the FHE-versus-MPC
crossover is a property of \emph{the baseline's preprocessing} rather than of
either cryptosystem, and disappears once MPC is given somewhat-homomorphic
preprocessing; active security on the nonlinear surface of private inference
costs $2.63\times$ bytes but only $1.06\times$ rounds. We also report our own
withdrawn claims---a ZKP backend that proved nothing, a bit-decomposition that
leaked its inputs, and a comparison-protocol crossover whose location we got wrong
three times before measuring it---because a study that reports only what worked is
not
reporting.
\end{abstract}

%==============================================================================
\section{The comparability problem}
\label{sec:intro}

Ask how fast private inference is and the literature answers in incompatible
units. Gazelle~\cite{Gazelle}, Delphi~\cite{Delphi} and Cheetah~\cite{Cheetah}
report two-party hybrid systems in which the server holds the model in
plaintext. Falcon~\cite{Falcon} reports three-party replicated sharing on a LAN.
CKKS~\cite{CKKS} papers report amortized throughput over batched slots. SNARK
papers report prover time for a circuit that a different paper would express
with a different constraint count. Each is internally rigorous. Together they do
not compose.

Four obstacles recur.

\begin{enumerate}
  \item \textbf{Different hardware.} A $5\times$ difference between two systems
    is unfalsifiable when one ran on a server-class Xeon and the other on a
    laptop.
  \item \textbf{Different accounting.} ``Communication'' may mean payload bytes,
    serialized bytes, or bytes per party. ``Rounds'' may mean protocol rounds,
    messages, or interactive calls. These differ by more than the effects being
    measured.
  \item \textbf{Different security models.} Semi-honest and malicious figures
    appear in adjacent columns of secondary literature without the distinction
    surviving the citation.
  \item \textbf{Different functions.} Two \ReLU{} costs may be for different bit
    widths, different fixed-point scales, or---in the worst case we found in our
    own work---subtly different functions that nobody checked against a common
    reference.
\end{enumerate}

A survey assembled from such numbers inherits all four. The alternative we
pursue is to re-measure the paradigms in one harness where the obstacles are
controlled by construction: the same machine, the same accounting code,
machine-checked security levels, and a bit-exactness assertion against a shared
plaintext reference before any cost is quoted.

The cost of this approach is equally clear and we state it up front. Our
implementations are research-quality Rust, not the tuned C++ of the systems we
discuss. \textbf{Absolute figures here are not competitive with optimized
implementations and should not be read as such.} What the single harness buys is
\emph{internally valid ratios}: comparisons in which the confounds above are held
fixed. Where we compare against published numbers we say so, and treat the
comparison as weaker evidence than our internal ones.

\paragraph{Contributions.}
\begin{itemize}
  \item A \textbf{systematization} of four paradigms organized by what they
    actually cost, with the cost dimensions (rounds, bytes, compute,
    preprocessing) separated rather than collapsed into one ``performance''
    number (\cref{sec:background}).
  \item A \textbf{single-artifact measurement study}: roughly \num{70000} lines
    of Rust in which one program runs on every paradigm, with a validated
    network model and uniform accounting (\cref{sec:artifact,sec:method}).
  \item \textbf{Cross-paradigm findings} not visible from any single paradigm's
    literature (\cref{sec:cross}), including two crossovers whose location
    depends on parameters the original claims left implicit---one of which we
    ourselves mislocated three times before measuring it
    (\cref{sec:crossover-nonlin}).\vspace{-1pt}
  \item A \textbf{methodology for honest benchmarking} (\cref{sec:method}), with
    the record of what it caught in our own work (\cref{sec:withdrawn}).
\end{itemize}

%==============================================================================
\section{Scope}
\label{sec:scope}

\paragraph{In scope.} Semi-honest and malicious three-party MPC over rings;
garbled circuits in the honest-majority setting; RLWE-based FHE (BFV~\cite{FV},
BGV~\cite{BGV}) and TFHE~\cite{TFHE}; FHE$\leftrightarrow$MPC
scheme switching; SNARKs over 64-bit words (Binius64~\cite{Binius}); and the
supporting primitives these need---oblivious transfer, private set intersection,
and Beaver-triple~\cite{Beaver} preprocessing from both OT and
somewhat-homomorphic encryption (SHE).

\paragraph{Out of scope.} Differential privacy and federated learning (a
different threat model: they protect an \emph{aggregate}, not a computation);
trusted execution environments (hardware trust assumption); ORAM (surveyed but
not implemented, see \cref{sec:open}); and transformer nonlinearities such as
softmax and LayerNorm, which need exponential, reciprocal and inverse-square-root
subprotocols we do not have.

\paragraph{Application focus.} We use private neural-network inference as the
running workload, because it is the one application that exercises every
paradigm: dense linear algebra (where FHE packing shines), comparisons (where
MPC and garbled circuits compete), and verifiable computation (where SNARKs
apply). Where a result is specific to this workload we say so.

%==============================================================================
\section{Four paradigms, four cost structures}
\label{sec:background}

The paradigms are usually taxonomized by security property. For a practitioner
choosing among them, the more useful axis is \emph{what each is expensive in}
(\cref{tab:paradigms}).

\begin{table}[t]
\centering
\small
\begin{tabular}{@{}lllll@{}}
\toprule
Paradigm & Secret from & Expensive in & Cheap in & Rounds \\
\midrule
Secret-sharing MPC & computing parties & \textbf{rounds}, bandwidth & local arithmetic & $O(\text{depth})$ \\
Garbled circuits   & computing parties & \textbf{bandwidth} & rounds & $O(1)$ \\
FHE                & computing server  & \textbf{compute}, ct.\ size & rounds (none) & none \\
ZK proofs          & the verifier      & \textbf{prover compute} & verifier, proof size & none \\
\bottomrule
\end{tabular}
\caption{The paradigms organized by cost structure rather than security
property. Three trade-offs follow directly, and drive most of this paper.}
\label{tab:paradigms}
\end{table}

\paragraph{MPC and garbled circuits trade rounds against bytes.}
This is the classic Yao~\cite{Yao}-versus-GMW trade. Neither dominates: the
winner depends on round-trip time and bandwidth, and the boundary is an
empirical question. \Cref{sec:crossover-nonlin} measures it, and reports the three
answers we published before the current one.

\paragraph{FHE and MPC trade compute against communication.}
FHE sends a lot once and computes hard; MPC sends little, many times. On a fast
LAN, MPC usually wins; on a thin pipe, the balance shifts. This framing omits
\emph{preprocessing}, and \cref{sec:crossover-fhe} shows the omission is what
makes most published FHE-versus-MPC crossovers fragile.

\paragraph{ZK is orthogonal; conflating it with the others is a category error.}
MPC and FHE keep data secret \emph{from the parties computing on it}. A SNARK
proves a computation was performed correctly while hiding the witness \emph{from
a verifier}---the prover holds everything in the clear. A deployment needing both
confidentiality and integrity pays for both. We keep them in separate tables for
this reason, and our benchmark prints the caveat alongside its numbers so they
cannot travel without it.

\subsection{The hybrid, and why it exists}

Neural-network inference alternates linear layers with nonlinearities. FHE is
good at the former and bad at the latter---comparison is not a low-degree
polynomial. MPC is the reverse: comparison costs a bit decomposition, but a
linear layer costs one round regardless of width.

Linear layers under FHE, nonlinearities under MPC, with a \emph{scheme switch}
between them, is therefore the natural design, and is what
Gazelle~\cite{Gazelle}, Delphi~\cite{Delphi} and Cheetah~\cite{Cheetah} all
instantiate. \Baru{} implements it and adds two things (\cref{sec:hybrid}): the
weights are encrypted too, and the switch is used as a \emph{noise refresh},
which removes bootstrapping from deep networks entirely.

%==============================================================================
\section{The \texorpdfstring{\protect\normalfont\scshape Baru}{Baru} artifact}
\label{sec:artifact}

\subsection{One program, many backends}

\Baru{} is a Lisp-like language whose interpreter manipulates secret values only
through opaque handles and a fixed set of primitive operations. This is the
Arithmetic Black Box abstraction of the SPDZ line~\cite{DamgardNielsen,SPDZ}: an
ideal functionality $\Fabb$ storing values by handle and exposing
\bt{add}/\bt{mul}/\bt{lt}/\bt{share\_public}/\bt{reveal}. A backend implements
those primitives; the evaluator never inspects a representation.

\begin{lstlisting}
; private ReLU on one neuron -- every value is secret-shared
(let ((w :sint8 -1) (x :sint8 3) (b :sint8 0))
  (let ((z (+ (* w x) b)))
    (reconstruct (* z (> z 0)))))          ; => 0

; a value computed under MPC crosses into FHE without being revealed
(let ((keys (fhe-keygen)) (x (share 6 32)) (y (share -7 32)))
  (fhe-dec keys (share->fhe keys (* x y))))   ; => -42
\end{lstlisting}

The consequence for a survey is direct: \emph{the same source text} is the
workload for every paradigm, so a cross-paradigm comparison cannot be confounded
by two implementations computing subtly different things. This is not
hypothetical; \cref{sec:withdrawn} records a case where it caught exactly that.

\subsection{Composability}

Because every gadget (\bt{relu}, \bt{max2}, \bt{maxpool2d}, \bt{argmax},
\bt{dense}, \bt{conv2d}) is a pure composition of $\Fabb$ primitives, universal
composability~\cite{Canetti} gives security of any network built from them for
free, given the primitives. Architecture search therefore needs no new security
argument per architecture, and the research question becomes ``which
architecture is \emph{fastest} under the cost model?'' rather than ``is this one
secure?''

\subsection{What is implemented}

\begin{table}[t]
\centering
\small
\begin{tabular}{@{}p{0.16\linewidth}p{0.78\linewidth}@{}}
\toprule
Layer & Contents \\
\midrule
MPC & ABY3~\cite{ABY3} replicated shares over $\Z_{2^{128}}$; \textbf{triple-free} multiplication from pairwise PRFs; bit decomposition (carry-save $3\to2$ plus Kogge--Stone~\cite{KoggeStone}); batched primitives; \textbf{malicious security with abort} ($\le 1$ of 3 corrupt), following Chida et al.~\cite{Chida} in the ring adaptation of Eerikson et al.~\cite{Eerikson}; in-process and TCP transports \\
Garbled & Half-gates~\cite{HalfGates}, free-XOR~\cite{FreeXOR}, fixed-key AES; comparator and fused $\ReLU\circ\text{rescale}$ circuits; OT-free label delivery in the three-party honest-majority setting \\
Falcon & $\Pi_{\mathrm{PC}}$ private compare over $\Z_p$, $\Pi_{\mathrm{WA}}$ \bt{wrap3}, $\Pi_{\mathrm{DReLU}}$, $\Pi_{\mathrm{SS}}$, $\Pi_{\mathrm{Maxpool}}$, exact truncation---a reproduction of Falcon~\cite{Falcon} inside the same harness \\
FHE & BFV~\cite{FV} and BGV~\cite{BGV} (depth-one, no evaluation keys), TFHE-rs~\cite{TFHErs,TFHE}, an exact real-arithmetic backend; multi-modular NTT with Shoup butterflies; CRT slot packing; Cheetah-style~\cite{Cheetah} convolution packing \\
Hybrid & MPC$\leftrightarrow$FHE switching both directions; re-encryption as a bootstrapping alternative; replicated threshold-key switch~\cite{AsharovMHE} \\
ZKP & Expression-to-circuit compiler; real Binius64~\cite{Binius} prover and verifier \\
Compiler & Backend-aware circuit IR, rewrite rules, cost models, protocol registry (dealer / Shamir / Semi2k~\cite{MPSPDZ} / LowGear~\cite{Overdrive}), per-operation protocol assignment \\
Primitives & Chou--Orlandi~\cite{ChouOrlandi} and IKNP~\cite{IKNP} OT, DH-PSI \\
\bottomrule
\end{tabular}
\caption{What the artifact implements. Every row is exercised by at least one
benchmark, and every engine attests its security level (\cref{sec:attest}).}
\label{tab:implemented}
\end{table}

\Cref{tab:implemented} summarizes. All figures below are from the maintainer's
machine (Apple M1 Pro, 8 cores), release build, unless noted.

%==============================================================================
\section{The paradigms, surveyed and measured}
\label{sec:paradigms}

\subsection{Secret-sharing MPC}
\label{sec:mpc}

\Baru{}'s MPC layer is ABY3-style~\cite{ABY3}: $x = s_0+s_1+s_2 \bmod 2^{128}$,
with party $i$ holding the replicated pair $(s_i, s_{i+1})$. Addition and scalar
multiplication are local. Multiplication is \textbf{triple-free}---each party
sends one ring element to one neighbour, with correlated randomness regenerated
locally from a one-time pairwise PRF key exchange. There is no dealer and no
preprocessing phase.

The engineering lesson this layer taught, four separate times, generalizes
beyond the artifact.

\finding{The arithmetic was never the bottleneck; the round structure was. Each
time a segment-level measurement said ``MPC is slow'', a per-phase breakdown
said ``one unbatched loop inside it is slow''.}

Batched forms cost \emph{the same number of rounds for any batch width}.
Adopting them moved the CNN's MPC segment from \SI{40.8}{\second} to
\SI{6.1}{\second} ($6.7\times$), with messages falling from \num{29211} to
\num{5591}. A reader should take from this that published MPC costs are highly
sensitive to batching discipline, and that comparing two systems' round counts
without checking whether both batch is comparing implementations, not protocols.

\paragraph{Cost of active security.}
Verified 2026-08-01 (\cref{tab:malicious}).

\begin{table}[t]
\centering
\small
\begin{tabular}{@{}lccc@{}}
\toprule
Workload & bytes & interactive calls & wall clock (WAN, \SI{25}{\milli\second}) \\
\midrule
\textbf{Comparisons}, $n{=}32$, 32-bit & $\mathbf{2.63\times}$ & $\mathbf{1.06\times}$ & $\mathbf{1.26\times}$ \\
Dot product, $n{=}32$, 32-bit & $2.47\times$ & $1.75\times$ & $2.80\times$ \\
\bottomrule
\end{tabular}
\caption{Semi-honest ABY3 versus malicious-with-abort, same workload, same
transport. Active security costs bytes, not rounds.}
\label{tab:malicious}
\end{table}

The comparison row is the one that matters for inference, since every \ReLU{},
max-pool and argmax reduces to a comparison.

\finding{Active security costs bytes, not rounds: the MAC doubles what is sent
but batches into the same round.}

The dot-product row is a small-workload artifact---the semi-honest leg needs only
four interactive calls, so a fixed pair-pool cost amortizes against nothing---and
we report it rather than hide it.

A second result generalizes, \textbf{with a correction we owe to
\cref{sec:crossover-nonlin}}. The statistical slack that ring-based verification
normally charges as a widening to $\Z_{2^{k+s}}$~\cite{SPDZ2k} is already
\emph{present} when the implementation shares over a ring much wider than its values:
Q16.16 values in a 128-bit ring leave \textbf{59 bits} of statistical security.

We previously added that this came at \emph{no widening cost}. That is false, and
packing the semi-honest paths is what exposed it: the 128-bit ring had been costing
$3.1$--$5.6\times$ in bytes on every one of them, which is exactly what narrowing to
the values' own width recovered. The verified path cannot recover it---its security is
$128 - \mathit{used\_bits}$, so narrowing to the value width sets $\lambda$ to zero.
The slack was never free; it was invisible, because nothing was packed and so nothing
revealed its price. The $2.63\times$ above is therefore a fair ratio at equal width
that \emph{understates} the cost of active security against a packed baseline, by a
factor bounded by the packing we measured on comparable paths. We have not run that
combination and claim no figure for it.

\subsection{Garbled circuits}
\label{sec:gc}

In the \emph{three-party honest-majority} setting, garbled-circuit input labels
are delivered by the co-garbler with no oblivious transfer. Since OT is the
expensive part of two-party garbling, this changes the economics substantially,
and is why garbling is competitive here at all.

Measured head-to-head against bit decomposition through the identical transport
(\num{1352} comparisons, 20-bit): \bt{a2b} needs 9 rounds and
\SI{1.29}{\mega\byte}; the garbled circuit needs 1 round and
\SI{1.69}{\mega\byte}. \textbf{Bit decomposition is therefore the byte-cheaper of the
two}, by $1.3\times$, and garbling's advantage is entirely in rounds: wall clock
favours it $18.8\times$ on a LAN and $4.0\times$ on an \SI{80}{\milli\second} /
\SI{100}{\mega\bit\per\second} WAN, the ratio shrinking as bandwidth rather than
latency binds.

\paragraph{This reverses a figure we published.} The same experiment previously
reported \bt{a2b} at \SI{5.94}{\mega\byte}, which made garbling look $3.5\times$
cheaper in bytes \emph{as well as} faster. That $4.6\times$ difference was \bt{u128}
serialisation padding: a 20-bit ring share travelled in 128 bits, and \bt{a2b} is
almost entirely ring traffic while garbled tables are already compact---so the
overhead fell on one side of the comparison only. With the encoding fixed
(\cref{sec:crossover-nonlin}) the byte ordering flips and only the round advantage
survives. Both routes remain bit-exact against the plaintext reference.

\paragraph{A correction to our own earlier claim.} The figures above are for the
\emph{garbling primitive in isolation}. The deployable three-party protocol,
including the $Y2A$ return leg that converts the boolean output back to an arithmetic
share, is smaller: $7.7\times$ on a LAN falling to $1.8\times$ at
\SI{80}{\milli\second} / \SI{100}{\mega\bit\per\second}, and $3.9$--$2.4\times$ for
the fused ReLU$\circ$rescale. We record both, because the gap between them is exactly
what gets lost in citation.

An earlier version of this paragraph quoted $1.4$--$1.7\times$ for the deployable
path, taken not from this experiment but from the CNN's \emph{MPC segment}---a
different benchmark on a different workload, presented as though it corrected this
one. That segment measures $6.1\times$ today. The old figure was right when written
and is stale for a reason we record in \cref{sec:withdrawn}: an unbatched
scheme-switch import was then 94\% of the MPC segment, loading both routes equally and
compressing the ratio toward~1. Both corrections illustrate one rule---quote the
deployable protocol, on the workload being claimed about, and name the benchmark that
produced the number.

\subsection{Fully homomorphic encryption}
\label{sec:fhe}

\Baru{} uses the cheapest FHE primitive that exists: one
ciphertext$\times$ciphertext multiplication, \emph{no relinearization keys, no
bootstrapping, no evaluation keys}, decrypting the degree-2 product directly
under $s^2$. Parameters are $n=4096$, $q=2^{100}$, $t=2^{20}$ for BFV---inside
the ${\approx}109$-bit $\log q$ ceiling for 128-bit security at that
dimension~\cite{HEStandard}.

Two packing results are worth surveying.

\begin{itemize}
  \item \textbf{Reversed-coefficient packing} makes a whole matrix--vector
    product one multiplication. We find this is \emph{already optimal for dense
    layers}: a degree-$n$ product admits at most $n/\mathrm{in}$ separable
    outputs, so $\lceil \mathrm{out}\cdot\mathrm{in}/n\rceil$ multiplications is
    the limit. Gazelle's SIMD packing~\cite{Gazelle} matches the throughput but
    requires rotations and Galois keys, forfeiting the no-evaluation-keys
    property.
  \item \textbf{Cheetah-style convolution packing}~\cite{Cheetah} makes a whole
    output channel one multiplication, independent of image size. Measured A/B at
    $28\times28$: \num{1352} multiplications become \textbf{2}, and FHE time falls
    from \SI{77.5}{\second} to \SI{264}{\milli\second}.
\end{itemize}

\finding{BGV is indistinguishable from BFV at identical parameters (repeated runs
put the gap at $+0.5\%$, $+0.5\%$, $-0.1\%$). Where the message sits in the
ciphertext, high bits or low, is not where the cost is.}

On arithmetic, our multi-modular NTT---NTT-friendly primes used as a
computational device and CRT-reconstructed, because a power-of-two $q$ admits no
direct NTT---is $1.8\times$ faster than NTL+GMP~\cite{NTL} at every size after
moving its butterfly to Shoup-preconditioned multiplication. NTL is retained as a
\emph{differential oracle} rather than a dependency, which is what made rewriting
the hottest code in the repository safe.

\subsection{The hybrid and the no-bootstrap result}
\label{sec:hybrid}

The construction's central observation is that the scheme switch performed for a
\ReLU{} \emph{also refreshes the ciphertext}. After switching back, each party
freshly encrypts its share component and the evaluator sums; the sum of fresh
encryptions is a fresh encryption. The FHE noise budget therefore resets at every
layer boundary, and a network of arbitrary depth runs with \textbf{zero FHE
bootstraps} at cost linear in depth, on an engine supporting only a
\emph{single} multiplication (\cref{tab:deep}).

\begin{table}[t]
\centering
\small
\begin{tabular}{@{}rccrrrcc@{}}
\toprule
depth & exact? & total & fhe & switch & mpc & our boots. & levelled-competitor boots. \\
\midrule
2  & \checkmark & \SI{145}{\milli\second}  & 20  & 59  & 54  & \textbf{0} & 0 \\
8  & \checkmark & \SI{539}{\milli\second}  & 81  & 239 & 210 & \textbf{0} & 3 \\
32 & \checkmark & \SI{2119}{\milli\second} & 333 & 958 & 820 & \textbf{0} & 15 \\
\bottomrule
\end{tabular}
\caption{Deep private inference, width 16, weights \emph{and} activations
encrypted, real ABY3 protocol, BFV at $n=4096,\ q=2^{100}$. Linear in depth and
bit-exact at every depth; the exactness is the evidence for the noise claim,
since accumulated noise would make deep outputs diverge.}
\label{tab:deep}
\end{table}

This construction also encrypts the \textbf{weights}. That is a stronger privacy
model than Gazelle, Delphi and Cheetah, whose server holds the model in
plaintext.

\subsection{Zero-knowledge proofs}
\label{sec:zkp}

Proving a Q16.16 quantized neuron, $\ReLU((\sum_i w_i x_i) \gg 16 + b)$, with
public weights and private inputs, on a real Binius64 prover
(\cref{tab:zkp}).

\begin{table}[t]
\centering
\small
\begin{tabular}{@{}rrrrr@{}}
\toprule
width & constraints & prove & verify & proof bytes \\
\midrule
1   & 14    & \SI{6.3}{\milli\second}  & \SI{1.0}{\milli\second} & \num{18544} \\
64  & 589   & \SI{8.0}{\milli\second}  & \SI{1.3}{\milli\second} & \num{47808} \\
\textbf{784} & \textbf{\num{7159}} & \textbf{\SI{32.3}{\milli\second}} & \textbf{\SI{4.8}{\milli\second}} & \textbf{\num{111712}} \\
\bottomrule
\end{tabular}
\caption{Constraints grow linearly ($\approx 9$ per input); proof size does
not---a $286\times$ larger circuit yields a $6\times$ larger proof. Width 784 is
the MNIST input dimension.}
\label{tab:zkp}
\end{table}

Two implementation findings generalize to anyone building arithmetic circuits
over 64-bit words. First, comparison gates are \textbf{unsigned}, so signed
comparison must flip sign bits first; without it $\ReLU(-1)$ returns $-1$---a
circuit perfectly sound about the wrong function. Second, and sharper:

\finding{Division cannot be checked the obvious way. A circuit cannot divide, so
the prover supplies quotient and remainder as witnesses; but checking $a = bq+r$
with $|r|<|b|$ is \emph{unsound}, because circuit multiplication wraps mod
$2^{64}$. For $a=1, b=3$ a cheating prover sets $r=0$ and $q=3^{-1} \bmod 2^{64}$
and both checks pass. Requiring the 128-bit product's high word to be the sign
extension of the low word closes it.}

We drive that forgery through the real prover as a test.

\subsection{Preprocessing: OT versus SHE}
\label{sec:triples}

Beaver-triple preprocessing is where the FHE-versus-MPC comparison is usually
decided, and it is usually decided implicitly (\cref{tab:triples}).

\begin{table}[t]
\centering
\small
\begin{tabular}{@{}llrrr@{}}
\toprule
route & domain & B/triple & triples/s & rounds \\
\midrule
LowGear (SHE)~\cite{Overdrive} & $\mathbb{F}_p$, 60-bit prime & \textbf{192} & \num{97304} & 2 \\
Semi2k (OT)~\cite{MPSPDZ}      & $\Z_{2^{64}}$ & \num{3074} & \num{2681} & 3 \\
\bottomrule
\end{tabular}
\caption{Two parties, semi-honest. \textbf{The two rows are over different
domains}; see the caveat below for the apples-to-apples figure.}
\label{tab:triples}
\end{table}

\paragraph{A comparability caveat the literature routinely drops---and which we
dropped ourselves.} A ${\approx}60$-bit prime field does \emph{not} represent
every 64-bit value, so \cref{tab:triples} compares a 61-bit field triple against
a 64-bit ring triple. At \textbf{equal 64-bit width} LowGear needs two CRT limbs,
giving \SI{384}{\byte}/triple against OT's \num{3074}---an $\mathbf{8.0\times}$
communication advantage, not the $16\times$ a naive reading suggests. Our own
experimental record quoted $16.3\times$; we correct it here. This is precisely
the error class a single-harness study is positioned to catch, and it is the
reason \cref{sec:method} exists.

The measured OT cost also validates the analytic model: \num{3074} measured
against \num{3136} predicted, a \SI{2}{\percent} error.

%==============================================================================
\section{Methodology}
\label{sec:method}

This is the section we would most like other artifacts to adopt.

\subsection{Identical transport, uniform accounting}

Every protocol---arithmetic MPC, garbled, Falcon, hybrid---runs over the same
transport stack, wrapped in counting and latency decorators that observe traffic
without touching protocol code. Both legs of any comparison pay real
serialization, real message passing and real byte counting. There is no
measurement asymmetry in which one side is simulated and the other is not.

This matters more than it sounds. The in-process simulator \emph{is not the
protocol}: its share multiplication reconstructs and reshares locally, which is
free and insecure, and it understated the pure-MPC baseline by ${\approx}2.7\times$
(\SI{6.7}{\milli\second} simulated against \SI{18}{\milli\second} real). Every
MPC number in this study is the real wire protocol, with three party threads
exchanging serialized messages.

\subsection{A validated network model}
\label{sec:netmodel}

The cost model charges $\text{rounds} \times \text{RTT} + \text{bytes} /
\text{bandwidth}$, letting a benchmark ``move'' parties between continents
without deploying anything. It is \textbf{validated against a real loopback TCP
socket}: worst error \SI{13.6}{\percent}, under \SI{1}{\percent} through the
middle of both sweeps, with the link's constants measured rather than assumed.
Cost really is linear in rounds and linear in bytes with no missing cross-term,
which is what licenses computing wall-clock from measured counts.

Building this validation found two harness bugs that had been silently distorting
results: per-message latency timers could \emph{reorder} back-to-back messages,
and the bandwidth model did not serialize bulk transfers cumulatively,
\emph{undercharging exactly the byte-heavy phases} the crossover experiments are
about. Both now have regression tests. We mention this because a cost model never
checked against a socket should be assumed to have bugs of this kind.

The honest limitation: the constants are loopback's (RTT \SI{23}{\micro\second}),
not a WAN's, so a wide-area deployment needs its own two constants.

\subsection{Bit-exactness as a gate on quoting cost}

\finding{A cost comparison between routes computing different things is
meaningless.}

Every backend is therefore checked against a common plaintext reference before
any number is reported, and the benchmarks refuse to present a clean table when a
cell disagrees. The encrypted CNN reproduces the plaintext prediction on 600 of
600 embedded medical images, run twice---before and after two optimizations---with
both passes agreeing on all 600. That is what makes those optimizations
\emph{exact} rather than merely close.

\subsection{Machine-checked security attestations}
\label{sec:attest}

Every engine attests the security level it actually implements, in one auditable
table, and asking a session for more than an engine attests is an \emph{error
rather than a silent downgrade}. Exactly one engine attests malicious security.
The editing rule is that an attestation records what the code \emph{does}, never
what the roadmap intends: raising one is the last step of an upgrade, after the
cheating-detection tests pass.

We consider this the most transferable idea in the artifact. Security level is
the dimension most often lost when numbers move between papers, and making it a
machine-checked property of running code---rather than prose in a limitations
section---is cheap to implement and hard to get wrong by accident.

\subsection{Recording withdrawn results---and deleting the component}
\label{sec:withdrawn-policy}

Negative, superseded and withdrawn results are kept in the experimental record
alongside the wins, with the reasoning that produced them.
\Cref{sec:withdrawn} summarizes what this caught.

The half that is easy to skip is the other one: \textbf{when a measurement
condemns a component, the code is removed and the number retained}. An artifact
that accumulates every paradigm it ever measured becomes unmaintainable and,
worse, misleading---a reader assumes everything present is recommended. Removed on
this basis: single-server PIR (loses to downloading the database by
${\approx}13\times$, with \emph{no} crossover in $N$); the SEAL/CKKS backends
(approximate comparison costs seconds per operation with real error, against
milliseconds and bit-exact under MPC, and carried \SI{3.4}{\mega\byte} of
vendored C++); an NTL polynomial backend (the hand-written NTT beat it
$1.8\times$ at every size); and an unsound ZKP placeholder
(\cref{sec:withdrawn}).

Two consequences. \textbf{Figures here for a removed component are not
reproducible from the current artifact}, and the experimental record says so at
the point of measurement. And where a removed measurement still feeds a live
claim, it survives as the smallest thing that can carry it: the
levelled-competitor bootstrap counts in \cref{tab:deep} come from the level
arithmetic of the CKKS sign approximation, now a 40-line dependency-free module
whose figures are pinned by a test---precisely because the apparatus that produced
them is gone.

%==============================================================================
\section{Cross-paradigm findings}
\label{sec:cross}

These are the results that motivated the single-harness approach, in the sense
that no single paradigm's literature could produce them.

\subsection{The FHE-versus-MPC crossover belongs to the baseline's preprocessing}
\label{sec:crossover-fhe}

The standard claim is that FHE overtakes MPC past some layer width, because MPC
preprocessing grows linearly while a packed ciphertext is flat. \Baru{}'s cost
model predicts it, and measurement confirms it---\emph{against an OT-preprocessed
baseline} (\cref{tab:layer}).

\begin{table}[t]
\centering
\small
\begin{tabular}{@{}rrrrr@{}}
\toprule
width & dealer (lower bound) & Semi2k (OT) & \textbf{LowGear (SHE)} & hybrid/BFV \\
\midrule
64   & \SI{600.6}{\milli\second} & \SI{694.5}{\milli\second} & \textbf{\SI{606.2}{\milli\second}} & \SI{692.1}{\milli\second} \\
1024 & \SI{603.0}{\milli\second} & \SI{942.8}{\milli\second} & \textbf{\SI{623.4}{\milli\second}} & \SI{692.1}{\milli\second} \\
4096 & \SI{610.9}{\milli\second} & \SI{1.74}{\second}        & \textbf{\SI{678.4}{\milli\second}} & \SI{733.7}{\milli\second} \\
\bottomrule
\end{tabular}
\caption{One private layer with both operands secret plus a \ReLU{}; WAN, two
parties, preprocessing charged per triple \emph{consumed}.}
\label{tab:layer}
\end{table}

Against \textbf{SHE-preprocessed} MPC the crossover is gone in this range
entirely. LowGear tracks the trusted-dealer lower bound within
\SI{5}{\milli\second} at width 64.

\finding{The honest statement is not ``the hybrid beats dishonest-majority MPC
past width 64'' but ``the hybrid beats \emph{OT-preprocessed} dishonest-majority
MPC past width 64, and loses to SHE-preprocessed MPC at every width measured.''
The mechanism survives---preprocessing volume decides it---but the boundary is a
property of the baseline, not of either cryptosystem.}

A second boundary moves with the \textbf{threat model}. Under an honest majority
(ABY3, multiplication from correlated randomness, zero preprocessing) the layer
stays on MPC at \emph{every} width. Under a dishonest majority with OT
preprocessing, FHE takes it past ${\approx}256$ multiplications. This explains a
structural fact about the literature that is otherwise puzzling: \emph{every
hybrid FHE+MPC inference system is two-party}. Gazelle, Delphi and Cheetah sit
precisely in the regime where the model says the hybrid pays.

\subsection{Protocol assignment can be derived, not annotated}

Given measured costs for each backend, a search over per-operation-kind
assignments returns $\bt{mul}\to\text{fhe}$, $\bt{cmp-lt}\to\text{binary}$ for a
256-wide encrypted layer with a comparison---that is, it \emph{derives} the
Gazelle/Delphi/Cheetah architecture rather than having it annotated. The prior
work on protocol assignment (HyCC~\cite{HyCC} $\to$ Silph~\cite{Silph} $\to$
SING~\cite{SING}) is \textbf{MPC-only}; extending the cost model to span MPC
\emph{and} FHE is what makes the hybrid architecture an output of the search
rather than an input.

Two structural non-linearities make this objective non-separable, which is worth
noting for anyone tempted to formulate it as an integer linear program: a
dependence level costs the \emph{maximum} over its instructions, not their sum;
and conversions \emph{batch per level}, so a crossing's cost depends on whether
another crossing already happens there.

\subsection{The nonlinearity crossover, and three wrong answers on the way to it}
\label{sec:crossover-nonlin}

\Baru{} implements three interchangeable comparison protocols---batched bit
decomposition, garbled circuits, and Falcon's \bt{wrap3}-based
$\Pi_{\mathrm{DReLU}}$~\cite{Falcon}---behind one interface, verified to compute the
same function bit-exactly. \Cref{tab:crossover} measures them with \emph{both}
protocols charged for preprocessing and neither paying for serialisation padding.

\begin{table}[t]
\centering
\small
\begin{tabular}{@{}llrrrrrr@{}}
\toprule
RTT & protocol & 1000 & 100 & 40 & 30 & 20 & 10 \si{\mega\bit\per\second} \\
\midrule
\multirow{2}{*}{\SI{0}{\milli\second}}
 & gc     & \textbf{566}  & \textbf{2928} & \textbf{6672} & \num{8228}  & \num{12294} & \num{23525} \\
 & falcon & \num{1950} & \num{5459} & \num{6824} & \textbf{7898} & \textbf{10714} & \textbf{19297} \\
\midrule
\multirow{2}{*}{\SI{20}{\milli\second}}
 & gc     & \textbf{2273} & \textbf{4450} & \textbf{7811} & \textbf{9713} & \textbf{13515} & \num{24996} \\
 & falcon & \num{6865} & \num{8087} & \num{10645} & \num{12153} & \num{15029} & \textbf{23386} \\
\bottomrule
\end{tabular}
\caption{Milliseconds per encrypted MedMNIST diagnosis; winner in bold. Crossover
between $30$ and \SI{40}{\mega\bit\per\second} at RTT 0 and between $10$ and
\SI{15}{\mega\bit\per\second} at \SI{20}{\milli\second}; higher round-trip time moves
it \emph{down}.}
\label{tab:crossover}
\end{table}

The mechanism is in the same run: Falcon sends \SI{3.5}{\mega\byte} against
\SI{10.1}{\mega\byte} ($2.88\times$ cheaper) in \num{288} messages against \num{61}
($4.7\times$ dearer), so the terms trade and the boundary is where the trade
balances.

\paragraph{We reported three other answers first, and the sequence is the more useful
result.} Each was reproducible, computed the correct diagnosis, and agreed with our
cost model; the first matched a break-even prediction of \num{17.1} against
$15$--$20$ measured.

\begin{enumerate}
  \item \textbf{$15$--$20$\,\si{\mega\bit\per\second}.} Falcon's
    $\Pi_{\mathrm{WA}}$ preprocessing was dealt from a shared seed, free precisely
    because it was not secret---so this compared Falcon's \emph{online} cost with
    garbled circuits' \emph{total}.
  \item \textbf{No crossover at any bandwidth.} With a real offline phase Falcon
    measured \SI{32.5}{\mega\byte} against \SI{10.1}{\mega\byte} and appeared dearer
    in both terms. \textbf{82\% of those bytes were padding}: \Baru{} serialised ring
    shares as \bt{u128} whatever their width.
  \item \textbf{A tie.} The grid stopped at \SI{10}{\mega\bit\per\second}, above the
    boundary.
\end{enumerate}

Answer 1 was wrong in \emph{both} directions at once: dealt preprocessing flattered
Falcon while padding penalised it. Correcting both moves the boundary \emph{up}, so
Falcon wins a wider region than we originally claimed, for entirely different
reasons---the errors had partly cancelled, which is why the figure looked sound.

\textbf{Neither bit-exactness gating nor cost-model agreement detects any of the
three.} Both establish that a computation is correct and that its cost is
explicable; neither asks whether two measurements are \emph{commensurable}. What
caught these was decomposing each total into the terms the model charges---phases,
messages, bytes, and the padded share of bytes---and checking that both protocols
were charged for the same things.

One rule we now apply generally: a spot check is not evidence about an encoding.
Narrowing the ring wire to a fixed 21 bits was exact for the $2^{20}$ dataset and
wrong for \emph{every} image of the $2^{22}$ one, because the safe width is a
property of the model rather than a constant.

\paragraph{What we can still say.} The architectural claim survives and is worth
separating from the performance one:

\finding{Three implementations of one comparison, bit-exact against each other
behind a single interface, make the protocol backing every \ReLU{}, max-pool and
argmax a one-line choice rather than a rewrite---which is what made this error
detectable rather than shippable.}

What is no longer supported is that the choice should be made \emph{per
deployment}: on this workload there is currently nothing to choose, since garbled
circuits dominate both alternatives everywhere we measured. And Falcon's own
claim---that its \emph{online} phase is cheap---remains untested here.
\bt{lt\_falcon\_batch} runs preprocessing inline because this backend has no
offline/online split, so every figure in \cref{tab:crossover} is a total. Building
that split is the open work, and it is the limitation that most constrains this
section's conclusions in both directions.

\paragraph{The scope of the withdrawal.} One bit decomposition per element is an
artefact of \emph{our} construction rather than a lower bound. Writing
$w_{01} = \mathrm{wrap2}(x_0,x_1)$ and $y = (x_0+x_1) \bmod L$, one gets
$\delta = w_{01} \oplus [\,y + x_2 \geq L\,]$ (verified exhaustively at
$\ell = 6$). Party~0 holds both $x_0$ and $x_1$, so it computes $w_{01}$ and $y$
locally and for free, leaving a single comparison between a value it knows and a
value parties~1 and~2 know---a millionaires' problem, not another
$\Pi_{\mathrm{WA}}$, and one that the same no-OT three-party garbling used for the
\bt{gc} route settles in one round and ${\approx}\ell$ AND gates. We have not
implemented it and claim no cost for it. It is enough to fix the scope: what we
withdraw is \emph{our} crossover, and the correct reading of
\cref{tab:crossover} is ``\Baru{}'s Falcon loses to \Baru{}'s garbled circuits once
preprocessing is charged''---not that Falcon's DReLU fails to repay its
preprocessing, which these measurements do not show and the construction above
suggests is false.

%==============================================================================
\section{Case study: private medical image diagnosis}
\label{sec:medmnist}

MedMNIST v2~\cite{MedMNIST} distributes twelve biomedical datasets at
$28\times28$---exactly MNIST's~\cite{MNIST} shape---so a private-inference
pipeline runs on them unchanged, and the artifact stops being about handwritten
digits (\cref{tab:medmnist}).

\begin{table}[t]
\centering
\small
\begin{tabular}{@{}lrr@{}}
\toprule
 & PneumoniaMNIST & OrganAMNIST \\
\midrule
task & chest X-ray, 2 classes & abdominal CT, 11 organs \\
integer accuracy, full test set & \SI{84.78}{\percent} (624) & \SI{81.18}{\percent} (\num{17778}) \\
\SI{95}{\percent} Wilson CI~\cite{Wilson} & $[81.7, 87.4]$ & $[80.6, 81.8]$ \\
balanced accuracy / AUC & \SI{80.64}{\percent} / \textbf{0.9314} & \SI{80.49}{\percent} / \textbf{0.9590} \\
sensitivity / specificity & \SI{97.18}{\percent} / \textbf{\SI{64.10}{\percent}} & --- \\
published ResNet-18 ($28\times28$) & 0.854 / 0.944 & 0.935 / 0.997 \\
majority-class baseline & \SI{62.50}{\percent} & \SI{18.48}{\percent} \\
\textbf{encrypted $=$ plaintext} & \textbf{300/300} & \textbf{300/300} \\
FHE bootstraps & 0 & 0 \\
online latency / diagnosis & \textbf{\SI{268}{\milli\second}} & \SI{291}{\milli\second} \\
\bottomrule
\end{tabular}
\caption{Private diagnosis on real biomedical images, image \emph{and} weights
encrypted, only the class index revealed.}
\label{tab:medmnist}
\end{table}

Three points a privacy venue should care about.

\paragraph{Privacy costs nothing in accuracy, and this is shown rather than
asserted.} The integer model matches its float parent to within 0.2 points, and
the encrypted pipeline reproduces the plaintext diagnosis bit-exactly on every
one of 600 images.

\paragraph{Accuracy alone misreads a medical classifier.} PneumoniaMNIST's test
set is ${\approx}60\%$ positive, and this model earns its \SI{84.78}{\percent}
largely by \emph{over-calling disease}: \SI{97.2}{\percent} sensitivity against
\SI{64.1}{\percent} specificity. Balanced accuracy is \SI{80.64}{\percent}. None
of this is visible while only accuracy is reported, and none of it is a
cryptographic result---but a paper claiming private diagnosis has to report it.

\paragraph{The accuracy ceiling is the packing, not the crypto.} We initially
attributed the model's size to the FHE noise budget. That was wrong. Three limits
bind and noise is the third: spatially, $3\times3$ valid convolution plus pooling
takes $28 \to 13 \to 5 \to 1$, admitting at most three convolutions; by
\emph{packing}, $2\cdot \mathrm{in}_c \cdot H \cdot W \lesssim n$ caps the first
convolution at 12 channels before the second's input passes $n=4096$. That is a
packing constraint that relaxes at $n=8192$, not a limit of the construction.
Depth, meanwhile, really is free.

%==============================================================================
\section{What the methodology caught in our own work}
\label{sec:withdrawn}

A study's credibility rests partly on what its authors were willing to withdraw.
Each item below was found by the discipline of \cref{sec:method} and is recorded
in the artifact rather than quietly fixed.

\paragraph{A ZKP backend that proved nothing.} Until 2026-07-31 the shipped
backend's \bt{verify} never checked that a witness satisfied any constraint---
anyone holding the public circuit could forge an accepting proof. It is retained,
unused, so a soundness test can keep demonstrating precisely what it failed to
do. \textbf{Every ZKP number predating that date is withdrawn.}

\paragraph{A privacy break in bit decomposition.} The previous \bt{a2b} masked
its input with a value derived from a \emph{public} lockstep counter, so the input
was recoverable from the transcript by any observer. Everything built on it---
\ReLU{}, max-pool, argmax---leaked its activations. The advertised ``no extra
communication per bit'' was a \emph{symptom}: the carry chain was free because
every operand was already public. Fixing it cost $4.4\times$ MPC time and
${\approx}130\times$ bandwidth, later recovered by triple-free multiplication and
batching.

\paragraph{An insecure FHE scheme, removed rather than caveated.} The fast engine
was a scale-invariant NTRU~\cite{LTV,YASHE} port. Analysis showed it insecure at
the parameters the construction needs even at depth
one~\cite{ABDucas,KirchnerFouque,DucasvanWoerden}: the large LTV secret raises
the fatigue point to ${\approx}2^{59}$, but a single multiply with usable $t$
forces $q \gtrsim 2^{84}$. It was deleted, not kept as a ``fast reference'';
every result transferred bit-exact and the FHE column got slower.

\paragraph{A selective-failure channel that passes every functional test.}
Attaching MACs to bit-decomposition leaves by multiplication is unsound: a
corrupt party injects $\delta$, the first downstream gate carries MAC error
$\delta y$ with $y$ a secret bit, so the run aborts when $y=1$ and completes when
$y=0$.

\finding{Deferring the abort does not close this. Batching hides \emph{when} an
abort happens, not \emph{whether} it happens; those are different channels and
only the first is closed by batching. We had asserted the opposite before
building it.}

\paragraph{A crossover withdrawn three times, for three different reasons.} The
location in \cref{sec:crossover-nonlin} has been wrong three ways: measured against
dealt (hence free) preprocessing, so it compared one protocol's online cost with
another's total; then measured with Falcon inflated $3.2\times$ by serialisation
padding, from which we wrongly inferred that no crossover could exist anywhere; then
reported as a tie because the grid stopped above the boundary. It is at
$2$--$4$\,\si{\mega\bit\per\second}. \textbf{All three passed our bit-exactness gate
and agreed with our cost model}, the first even matching its break-even prediction.
A related lesson we now apply: a two-image spot check is not evidence about a wire
format---narrowing to 21 bits was exact for the $2^{20}$ dataset and wrong for
\emph{every} image of the $2^{22}$ one.

\paragraph{Comparison errors of framing.} Pricing FHE's \emph{linear} part
against MPC's \emph{whole layer} (omitting the \ReLU{} that FHE cannot do at
all); computing a crossover against a trusted-dealer baseline (asking whether FHE
beats free preprocessing); and the cross-domain triple comparison corrected in
\cref{sec:triples}.

\paragraph{Honest negatives retained.} PIR loses to downloading the database by
${\approx}13\times$ with no crossover in $N$. Hybrid-TFHE does not pay. Faster
ring arithmetic does \emph{not} bring FHE close to MPC. Per-layer bit-width
tuning was measured at ${\approx}10\%$ and not built. A fused garbled max-of-four
was measured as a $1.68\times$ \emph{loss} and not built.

%==============================================================================
\section{Threats to validity}
\label{sec:threats}

\paragraph{Implementation quality.} Research Rust, not tuned C++. Absolute
figures are not competitive; ratios measured within the harness are the
defensible output. Where we compare against published numbers---for example our
sign extraction at ${\approx}\SI{0.043}{\milli\second}$ against Falcon's
optimized C++---the gap is implementation, not protocol, and we say so.

\paragraph{Single machine.} All figures from one Apple M1 Pro. Cross-machine
variance is uncharacterized. The network model is validated on loopback, so WAN
constants are modelled.

\paragraph{The simulated network is a model.} Validated to \SI{13.6}{\percent}
worst case, but it is still $\text{rounds}\times\text{RTT} +
\text{bytes}/\text{bandwidth}$, not a real WAN with congestion, jitter or packet
loss.

\paragraph{One disclosed cryptographic gap.} FHE threshold \emph{key generation}
is dealer-based. It is recorded on the code and in the attestation table rather
than only in prose.

The Falcon backend's preprocessing used to be the second such gap---dealt from a
common seed, so all three parties could reconstruct the mask. It is now generated
by bit composition, at a cost of one bit decomposition per element offline. That
cost bears on how the crossover should be read: Falcon is cheaper than batched bit
decomposition \emph{online}, which is what \cref{tab:crossover} measures, and not
cheaper once the offline phase is added in.

\paragraph{Workload breadth.} Findings are anchored on feed-forward CNN inference
at $28\times28$. Transformers, larger images, and non-inference workloads are not
covered.

\paragraph{Not an adversarial evaluation.} Cheating-detection tests drive real
attacks over the transport for the malicious ABY3 layer, but this is not a
security audit, and no formal proofs are machine-checked.

%==============================================================================
\section{Open problems}
\label{sec:open}

\begin{enumerate}
  \item \textbf{A CIFAR-scale private CNN} needs SIMD-slot packing; coefficient
    packing bounds image size. The composition and security story is independent
    of it.
  \item \textbf{Transformer nonlinearities} (softmax, LayerNorm) need
    exponential, reciprocal and inverse-square-root subprotocols.
  \item \textbf{Malicious security at the FHE boundary.} Replication cannot catch
    a \emph{malformed input ciphertext}---both honest co-holders compute on it
    faithfully and agree. This sharpens the positioning: \emph{ZK is needed only
    at the input boundary, never for the evaluation}, where single-server
    verifiable FHE needs it over the whole computation. The boundary proof is not
    implemented.
  \item \textbf{Secret-index array access}, priced by the cost model, before
    ORAM. CirC~\cite{CirC} implements exactly the two cases either side of that
    line and does not implement ORAM at all.
  \item \textbf{Per-instruction protocol assignment}, versus the current
    per-operation-kind granularity.
  \item \textbf{Cross-machine and real-WAN validation} of the cost model's
    constants.
\end{enumerate}

%==============================================================================
\section{Conclusion}

Cross-paradigm performance claims in privacy-preserving computation are not
currently falsifiable, because the numbers come from incommensurable sources. We
have described a harness in which they are, and reported what changes when the
paradigms are re-measured under uniform accounting: an FHE-versus-MPC crossover
that belongs to the baseline's preprocessing rather than to either cryptosystem,
active security that costs bytes but not rounds; and a comparison-protocol boundary
at $30$--$40$\,\si{\mega\bit\per\second} that we mislocated three times first, each
time with a figure that was reproducible and agreed with our own cost model.
\Cref{sec:crossover-nonlin} reports that sequence rather than only its endpoint,
because the failure mode it illustrates---two measurements that are individually
sound and mutually incommensurable---is not caught by any check we had. Six claims
did not survive re-measurement, five of them our own.

The transferable part is not the artifact but the discipline: identical
transports for both legs of any comparison, a cost model checked against a real
socket, bit-exactness as a precondition for quoting cost, and security level as a
machine-checked property of running code rather than prose. Each is cheap to
adopt, and each caught something here that would otherwise have been published.

%==============================================================================
\section*{Reproducibility}

Single Rust crate, no GPU. \bt{cargo test} covers correctness; every
\bt{bench-*} subcommand emits \bt{--json}. All numbers are produced by these
commands; none are hand-entered.

\begin{center}
\small
\begin{tabular}{@{}ll@{}}
\toprule
Result & Command \\
\midrule
Deep inference, zero bootstraps & \bt{bench-deep-nn --depths 2,4,8,16,32} \\
Private diagnosis (\cref{tab:medmnist}) & \bt{bench-medmnist --dataset pneumonia} \\
Nonlinearity comparison (\cref{tab:crossover}, withdrawn) & \bt{bench-medmnist --sweep} \\
Cost of active security (\cref{tab:malicious}) & \bt{bench-malicious --n 32 --bits 32} \\
Triples: SHE vs OT (\cref{tab:triples}) & \bt{bench-triples} \\
FHE$\leftrightarrow$MPC crossover (\cref{tab:layer}) & \bt{bench-layer} \\
Cost model vs real socket (\cref{sec:netmodel}) & \bt{bench-net} \\
Proving a quantized neuron (\cref{tab:zkp}) & \bt{bench-zkp-layer --widths 1,...,784} \\
Falcon vs \bt{a2b} vs garbled & \bt{bench-falcon} \\
\bottomrule
\end{tabular}
\end{center}

\bibliographystyle{plain}
\bibliography{references}

\end{document}
